% Generated draft: requires mathematical and source review.
\documentclass[11pt]{article}
\usepackage[margin=25mm]{geometry}
\usepackage{fontspec}
\setmainfont{Latin Modern Roman}
\usepackage{amsmath,amssymb,mathtools,unicode-math}
\setmathfont{Latin Modern Math}
\usepackage{enumitem,needspace,xurl,hyperref}
\setlength{\parindent}{0pt}
\setlength{\parskip}{5pt}
\setlength{\emergencystretch}{3em}
\newcommand{\N}{\mathbb{N}}
\newcommand{\Z}{\mathbb{Z}}
\newcommand{\Q}{\mathbb{Q}}
\newcommand{\R}{\mathbb{R}}
\newcommand{\C}{\mathbb{C}}
\newcommand{\F}{\mathbb{F}}
\newcommand{\A}{\mathbb{A}}
\newcommand{\PP}{\mathbb{P}}
\newcommand{\E}{\mathbb{E}}
\newcommand{\eps}{\varepsilon}
\newcommand{\dd}{\,\mathrm{d}}
\begin{document}
\section*{UnsolvedMath ID: 700862}
\textbf{Source problem code:} catalogue entry
\subsection{Definitions and mathematical statement}
Let \(p\) be a prime. An abstract group \(G\) is a \(p\)-group if every element has finite order equal to a power of \(p\). It is finitely generated if \(G=\langle S\rangle\) for some finite subset \(S\subseteq G\). It is residually finite if, for every \(g\in G\setminus\{1\}\), there are a finite group \(Q\) and a homomorphism \(\varphi:G\to Q\) such that \(\varphi(g)\ne 1\). For a subgroup \(H\le G\), write \([G:H]\) for its index.

The supplied source describes the following as a well-known problem: does there exist an infinite, finitely generated, residually finite \(p\)-group \(G\) such that
\[
\forall H\le G,\qquad |H|<\infty\ \text{or}\ [G:H]<\infty?
\]
The quantifier ranges over all subgroups, without a finite-generation or normality restriction. The source does not specify whether \(p\) is prescribed, whether existence is required for every prime, or whether existence for some prime suffices; this qualification remains unresolved.
\subsection{Short English statement}
The question asks for an infinite, finitely generated, residually finite group whose elements all have orders that are powers of one prime, and whose every infinite subgroup has finite index. The intended quantification over the prime is unspecified.
\subsection{Sources}
\begin{itemize}
\item Local UnsolvedMath registry: permanent ID \(700862\), exact source problem code \texttt{catalogue entry}, title \emph{Kourovka Notebook Problem 12.43}. The supplied record labels the problem open; this label has not been independently verified.
\item \emph{The Kourovka Notebook (Unsolved Problems in Group Theory)}, problem 12.43, 12th issue (1992); also identified in the supplied record with the 21st issue edition (2026), arXiv:\texttt{1401.0300v48}. Proposed by J. C. Lennox and described in the supplied statement as a well-known problem. The supplied collection description names E. I. Khukhro and V. D. Mazurov as editors.
\end{itemize}
\end{document}
